\section{Formal Problem Statement}
\label{sec:problem}

We establish the formal foundation for the CausalOCPM framework, beginning
with the data model and concluding with a precise statement of the causal
policy evaluation problem.

\begin{definition}[Object-Centric Event Log]
\label{def:ocel}
Let $\mathcal{L} = (E, \mathcal{O}, \mathit{OT}, \pi_{\mathit{act}},
\pi_{\mathit{time}}, \pi_{\mathit{omap}})$ be an OCEL 2.0 event log where:
\begin{itemize}
    \item $E$ is a finite set of events;
    \item $\mathcal{O}$ is a finite set of objects;
    \item $\mathit{OT}$ is a finite set of object types;
    \item $\pi_{\mathit{act}}: E \to \mathcal{A}$ maps events to activities
          from a finite activity set $\mathcal{A}$;
    \item $\pi_{\mathit{time}}: E \to \mathbb{R}$ maps events to timestamps;
    \item $\pi_{\mathit{omap}}: E \to \mathcal{P}(\mathcal{O})$ maps each
          event to the set of objects it involves.
\end{itemize}
\end{definition}

\noindent The key distinction from classical (case-centric) event logs is that
each event may involve objects of multiple types simultaneously --- for instance,
a single \texttt{GoodsReceipt} event may reference a purchase order, a vendor,
and a material item at once. This richer structure is essential for capturing
cross-entity causal dependencies in real processes.

\begin{definition}[Object-Centric Causal Process Model (OCPM)]
\label{def:ocpm}
Given an OCEL event log $\mathcal{L}$, let
$\mathcal{V} = \{V_1, \ldots, V_k\}$ be a set of process variables derived
from $\mathcal{L}$ (e.g., via case-level aggregation).
An \emph{Object-Centric Causal Process Model} is a tuple
$\mathcal{M} = (\mathcal{V}, \mathcal{D}, \mathcal{F}, \mathcal{U})$ where:
\begin{itemize}
    \item $\mathcal{D} = (\mathcal{V}, \mathcal{E_D})$ is a directed acyclic
          graph (DAG) over $\mathcal{V}$;
    \item $\mathcal{F} = \{f_i : \mathit{Pa}(V_i) \times U_i \to V_i\}_{i=1}^{k}$
          are \emph{structural equations}, where $\mathit{Pa}(V_i)$ denotes
          the parents of $V_i$ in $\mathcal{D}$ and $U_i$ is an exogenous
          noise term;
    \item $\mathcal{U} = \{U_1, \ldots, U_k\}$ are mutually independent
          exogenous noise variables.
\end{itemize}
\end{definition}

\noindent The structural equations $\mathcal{F}$ encode the data-generating
process for each variable conditional on its causal parents. Unlike correlational
regression models, these equations support the \emph{do}-operator: setting
$V_i = v$ (an intervention) modifies $f_i$ but leaves all other equations intact,
enabling counterfactual reasoning over the process.

\begin{assumption}[Markov Condition]
\label{ass:markov}
Every variable $V_i \in \mathcal{V}$ is conditionally independent of its
non-descendants in $\mathcal{D}$ given its parents $\mathit{Pa}(V_i)$.
\end{assumption}

\begin{assumption}[Faithfulness]
\label{ass:faithfulness}
All conditional independencies in the observed joint distribution
$P(\mathcal{V})$ are entailed by the DAG $\mathcal{D}$ via the Markov
condition (Assumption~\ref{ass:markov}). Equivalently, no conditional
independence holds in $P(\mathcal{V})$ unless it is implied by the
$d$-separation relations in $\mathcal{D}$.
\end{assumption}

\begin{assumption}[Causal Sufficiency]
\label{ass:sufficiency}
There are no unmeasured common causes of any two variables in $\mathcal{V}$.
That is, all common causes of observed variables are themselves observed.
\end{assumption}

\noindent \textbf{Remark.} Assumption~\ref{ass:sufficiency} is the most
restrictive. In practice, processes always contain unmeasured variables.
We relax this assumption empirically in Section~\ref{sec:sensitivity} via
sensitivity analysis for unmeasured confounding, following the methodology
of \citet{cinelli2020making}. The sensitivity analysis quantifies how strong
an unmeasured confounder would need to be to invalidate our conclusions.

\begin{problem}[Causal Policy Evaluation in Object-Centric Process Mining]
\label{prob:main}
Given an OCEL event log $\mathcal{L}$ and a set of process variables
$\mathcal{V}$ derived from $\mathcal{L}$, and given a treatment variable
$T \in \mathcal{V}$ and outcome variable $Y \in \mathcal{V}$:

\textbf{Find:} A DAG $\mathcal{D}$ over $\mathcal{V}$ and a structural
causal model $\mathcal{M}$ such that for any intervention $do(T = t)$,
the interventional distribution $P(Y \mid do(T = t))$ is identified and
can be estimated from the observed data in $\mathcal{L}$.

\textbf{Subject to:}
\begin{enumerate}
    \item Assumptions~\ref{ass:markov}--\ref{ass:sufficiency} (with robustness
          to violation of~\ref{ass:sufficiency} quantified via sensitivity analysis);
    \item The OCEL multi-object event structure of $\mathcal{L}$, under which
          variables $T$ and $Y$ may derive from events referencing different
          object types in $\mathit{OT}$;
    \item The constraint that $\mathcal{D}$ may require cross-object-type
          causal edges (e.g., $\texttt{order\_complexity} \to
          \texttt{supplier\_a} \to \texttt{material\_lead\_time}$), where
          each variable in the path is associated with a distinct object type
          in $\pi_{\mathit{omap}}$.
\end{enumerate}
\end{problem}

\noindent The primary challenge in Problem~\ref{prob:main} is that standard
causal inference methods assume a flat, case-centric data model. We show in
Section~\ref{sec:theorem} that backdoor adjustment remains valid under the
OCEL multi-object structure, provided that the OCEL-to-variable mapping
preserves the conditional independence structure of the underlying causal DAG.
